\subsection{Construction of the Haar measure of a group from the Haar measures of certain subgroups}

Let \(\mathbf{G}\) be a locally compact group, \(\mathbf{X}\) and \(\mathbf{Y}\) two closed subgroups of \(\mathbf{G}\) such that \(\Omega = \mathbf{XY}\) contains a neighborhood \(\mathbf{U}\) of \(e\) . Then \(\Omega\) is open in \(\mathbf{G}\) ; for, for any \(x_0 \in \mathbf{X}\) and \(y_0 \in \mathbf{Y}\) , \(\mathbf{XY} = (x_0\mathbf{X})(\mathbf{Y}y_0) \supset x_0\mathbf{U}y_0\) , and \(x_{0}Uy_{0}\) is a neighborhood of \(x_{0}y_{0}\) ; thus \(\Omega\) is a neighborhood of each of its points.

*When G is a Lie group with Lie algebra g, the condition imposed on X and Y is satisfied if the subalgebras corresponding to X and Y have sum g .*

The group \(X \times Y\) operates continuously in G on the left, by the law \((x, y) \cdot s = xsy^{-1}\) ( \(x \in X, y \in Y, s \in G\) ). Let \(Z = X \cap Y\) . The stabilizer of e in \(X \times Y\) is the subgroup \(Z_{0}\) of \(X \times Y\) formed by the pairs \((z, z)\) , where \(z \in Z\) , a subgroup canonically isomorphic to Z. Thus the set \(\Omega\) may be identified with the homogeneous space \((X \times Y)/Z_{0}\) of left cosets; more precisely, the mapping \((x, y) \mapsto xy^{-1}\) of \(X \times Y\) onto \(\Omega\) defines, by passage to the quotient, a continuous bijection of \((X \times Y)/Z_{0}\) onto \(\Omega\) . We shall assume that this mapping is a homeomorphism. (This is notably the case if G is countable at infinity: cf.~App. I.)

PROPOSITION 13. --- Suppose in addition that Z is compact. Let \(\mu_{\mathrm{G}}\) , \(\mu_{\mathrm{X}}\) , \(\mu_{\mathrm{Y}}\) be left Haar measures on G, X, Y, and \(\Lambda\) the restriction of \(\Delta_{\mathrm{G}}\) to Y. Then the restriction \(\mu\) of \(\mu_{\mathrm{G}}\) to \(\Omega\) is, up to a constant factor, the image of \(\mu_{\mathrm{X}} \otimes (\Lambda^{-1} \cdot \mu_{\mathrm{Y}})\) under the mapping \((x, y) \mapsto xy^{-1}\) of X × Y onto \(\Omega\) (a mapping that is proper).

For \(x\in \mathbf{X},y\in \mathbf{Y}\)

\[
\boldsymbol {\gamma} \big ((x, y) \big) \mu = \boldsymbol {\delta} (y) \boldsymbol {\gamma} (x) \mu = \Delta_ {G} (y) \mu .
\]

Identifying \(\Omega\) with the homogeneous space \((\mathrm{X} \times \mathrm{Y}) / \mathrm{Z}_0\) and choosing a suitable Haar measure on \(\mathbf{Z}_0\) , one sees that \(\mu^\sharp\) is the product of the left Haar measure of \(\mathrm{X} \times \mathrm{Y}\) , namely \(\mu_{\mathrm{X}} \otimes \mu_{\mathrm{Y}}\) , by the function \((x, y) \mapsto \Delta_{\mathrm{G}}(y)^{-1}\) (No.~6, Lemma 6). On the other hand \(\mu\) is, up to a constant factor, the image of \(\mu^\sharp\) under the canonical mapping of \(\mathrm{X} \times \mathrm{Y}\) onto \(\Omega\) (No.~3, Remark 2).

COROLLARY. --- Let f be a function defined on \(\Omega\) , with values in a Banach space or in \(\overline{R}\) . For f to be \(\mu\) -integrable, it is necessary and sufficient that the function \((x,y)\mapsto\mathbf{f}(xy)\Delta_{\mathrm{G}}(y)\Delta_{\mathrm{Y}}(y)^{-1}\) be \((\mu_{\mathrm{X}}\otimes\mu_{\mathrm{Y}})\) -integrable, in which case

\[
\int_ {\Omega} \mathbf {f} (\omega) d \mu (\omega) = a \iint_ {\mathrm{X} \times \mathrm{Y}} \mathbf {f} (x y) \Delta_ {\mathrm{G}} (y) \Delta_ {\mathrm{Y}} (y) ^ {- 1} d \mu_ {\mathrm{X}} (x) d \mu_ {\mathrm{Y}} (y),\tag{23}
\]

where \(a\) is a constant \(>0\) independent of \(\mathbf{f}\) .

By Prop. 13, and Ch. V, §4, No.~4, Th. 2, for f to be \(\mu\) -integrable, it is necessary and sufficient that the function \((x,y)\mapsto\mathbf{f}(xy^{-1})\) be integrable for \(\mu_{\mathrm{X}}\otimes(\Lambda^{-1}\cdot\mu_{\mathrm{Y}})\) , or again that the function \((x,y)\mapsto\mathbf{f}(xy^{-1})\Delta_{\mathrm{G}}(y)^{-1}\) be integrable for \(\mu_{X}\otimes\mu_{Y}\) , or again that the function \((x,y)\mapsto\mathbf{f}(xy)\Delta_{\mathrm{G}}(y)\Delta_{\mathrm{Y}}(y)^{-1}\) be integrable for \(\mu_{X}\otimes\mu_{Y}\) . Formula (23) results from an analogous argument.

PROPOSITION 14. --- Suppose that the conditions of Prop. 13 are fulfilled and that, in addition, Y is normal.

\begin{enumerate}
\def\labelenumi{\alph{enumi})}
\item
  The restriction of \(\mu_{\mathrm{G}}\) to \(\Omega\) is, up to a constant factor, the image of \(\mu_{\mathrm{X}} \otimes \mu_{\mathrm{Y}}\) under the mapping \((x,y) \mapsto xy\) of \(\mathrm{X} \times \mathrm{Y}\) onto \(\Omega\) .
\item
  For \(x \in \mathbf{X}\) and \(y \in \mathbf{Y}\) ,
\end{enumerate}

\[
\Delta_ {\mathrm{G}} (x y) = \Delta_ {\mathrm{X}} (x) \Delta_ {\mathrm{Y}} (y) \bmod (i _ {x}),
\]

where \(i_x\) denotes the automorphism \(v \mapsto x^{-1}vx\) of Y.

We have \(\Delta_{\mathrm{G}} = \Delta_{\mathrm{Y}}\) on Y (Prop. 10 b)), therefore a) follows from (23). Let \(x_0 \in \mathbf{X}\) , \(y_0 \in \mathbf{Y}\) . Denote by \(p\) the mapping \((x, y) \mapsto xy\) of \(\mathbf{X} \times \mathbf{Y}\) onto \(\Omega\) . Since

\[
x y \left(x _ {0} y _ {0}\right) ^ {- 1} = x x _ {0} ^ {- 1} \left(x _ {0} y y _ {0} ^ {- 1} x _ {0} ^ {- 1}\right) = x x _ {0} ^ {- 1} i _ {x _ {0} ^ {- 1}} \left(y y _ {0} ^ {- 1}\right),
\]

we have

\[
\begin{array}{r l} \Delta_ {\mathrm{G}} (x _ {0} y _ {0}) p (\mu_ {\mathrm{X}} \otimes \mu_ {\mathrm{Y}}) & = \pmb {\delta} (x _ {0} y _ {0}) p (\mu_ {\mathrm{X}} \otimes \mu_ {\mathrm{Y}}) \\ & = p \big (\pmb {\delta} (x _ {0}) \mu_ {\mathrm{X}} \otimes i _ {x _ {0} ^ {- 1}} \pmb {\delta} (y _ {0}) \mu_ {\mathrm{Y}} \big) \\ & = p \big (\Delta_ {\mathrm{X}} (x _ {0}) \mu_ {\mathrm{X}} \otimes \Delta_ {\mathrm{Y}} (y _ {0}) (\bmod i _ {x _ {0}}) \mu_ {\mathrm{Y}} \big) \\ & = \Delta_ {\mathrm{X}} (x _ {0}) \Delta_ {\mathrm{Y}} (y _ {0}) (\bmod i _ {x _ {0}}) p (\mu_ {\mathrm{X}} \otimes \mu_ {\mathrm{Y}}), \end{array}
\]

whence b).

Remark. --- Prop. 14 applies in particular when G is the topological semi-direct product of X by Y (GT, III, §2, No.~10). In this case, \(Z = \{e\}\) and \(\Omega = G\) . Since \(yx = x i_{x}(y)\) for \(x \in X\) , \(y \in Y\) , the measure \(\mu_{G}\) is also, up to a constant factor, the image of \((\text{mod } i_{x})\mu_{X} \otimes \mu_{Y}\) under the mapping \((x, y) \mapsto yx\) of \(X \times Y\) into G.

